\section{Component-to-Module Integration Matrix}

\subsection{Status Codes}

The following codes are used in the integration tables.
\begin{itemize}
  \item \statuscode{D}: direct model consumption in the target module,
  \item \statuscode{F}: fixed injection or fixed parameter only,
  \item \statuscode{P}: used after projection into a canonical equivalent,
  \item \statuscode{N}: not currently consumed by that module.
\end{itemize}

\subsection{Module Preprocessing Differences}

The most important architectural difference between modules is whether they
reuse the canonical-projection path.
\begin{longtable}{L{0.20\linewidth}L{0.27\linewidth}L{0.43\linewidth}}
\toprule
Module family & Uses \texttt{project\_to\_canonical\_models}? & Practical consequence \\
\midrule
\endhead
Main PF family & yes & Flexible loads, asymmetric loads, transformers, switches, chargers, and optional three-phase models are converted before assembly. \\
Native / parity OPF & yes & Same canonical equivalents as PF are available to the OPF balances. \\
Short circuit & yes & Projected branches and loads are visible to the fault-network builder. \\
Distribution PF & no & Works directly on \texttt{ACSystem} / \texttt{HybridPowerSystem}; transformer, switch, flexible-load, asymmetric-load, and charger projection is not automatic. \\
ONR & no & Works directly on \texttt{ACSystem}; hybrid-only components and projection-only equivalents are not inserted automatically. \\
\bottomrule
\end{longtable}

\subsection{Detailed Integration Matrix}

For quick reading, the table legend is repeated here:
\statuscode{D} = direct,\;
\statuscode{F} = fixed-only,\;
\statuscode{P} = projection-based,\;
\statuscode{N} = not consumed.

\small
\begin{longtable}{L{0.34\linewidth}cccccccc}
\caption{Current component integration across solver modules}\label{tab:integration-current}\\
\toprule
Component model & \pfmark & \opfmark{} fixed & \opfmark{} var & \statuscode{DCOPF} & \statuscode{RPO} & \scmark & \dpfmark & \onrmark \\
\midrule
\endfirsthead
\toprule
Component model & \pfmark & \opfmark{} fixed & \opfmark{} var & \statuscode{DCOPF} & \statuscode{RPO} & \scmark & \dpfmark & \onrmark \\
\midrule
\endhead
\bottomrule
\endfoot
\texttt{ACBus} & D & D & N & D & D & D & D & D \\
\texttt{ACBranch} & D & D & N & D & D & D & D & D \\
\texttt{Transformer2W} & P & P & N & N & P & P & N & N \\
\texttt{Transformer3W} & P & P & N & N & P & P & N & N \\
\texttt{Switch} & P & P & N & N & P & P & N & N \\
\texttt{ExternalGrid} & D & F & N & N & F & D & N & N \\
\texttt{Generator} & D & D & D & D & D & D & D & F \\
\texttt{StaticGenerator} & D & D & N & F & D & N & D & F \\
\texttt{RenewableGen} & D & D & D & F & D & N & D & F \\
\texttt{PVSystem} & D & D & D & N & D & N & D & F \\
\texttt{Storage} (AC) & D & D & D & N & N & N & D & F \\
\texttt{Load} & D & D & F & D & D & D & D & F \\
\texttt{FlexibleLoad} & P & P & D & N & N & P & N & N \\
\texttt{AsymmetricLoad} & P & P & N & N & P & P & N & N \\
\texttt{Shunt} & D & D & N & N & D & F & F & N \\
\texttt{ChargingStation} & D & D & N & N & N & N & D & F \\
\texttt{Charger} & P & P & N & N & N & N & N & N \\
\texttt{AsynchronousMotor} & N & N & N & N & N & D & N & N \\
\texttt{DCBus} & D & D & N & N & N & N & N & N \\
\texttt{DCBranch} & D & D & N & N & N & N & N & N \\
\texttt{DCLoad} & D & D & N & N & N & N & N & N \\
\texttt{Storage} (DC) & D & D & D & N & N & N & N & N \\
\texttt{StaticGenerator} (DC) & D & D & N & N & N & N & N & N \\
\texttt{VSCConverter} & D & D & D & N & N & N & N & N \\
\texttt{DCDCConverter} & D & D & D & N & N & N & N & N \\
\texttt{EnergyRouterPort} & D & D & N & N & N & N & N & N \\
\texttt{EnergyRouter} & D & D & N & N & N & N & N & N \\
\texttt{VirtualPowerPlant} & D & D & N & N & N & N & D & N \\
\texttt{Microgrid} & D & D & N & N & N & N & D & N \\
\texttt{MobileStorage} & D & D & N & N & N & N & D & N \\
\texttt{ThreePhaseACBus} & P & P & N & N & P & P & N & N \\
\texttt{ThreePhaseACLine} & P & P & N & N & P & P & N & N \\
\texttt{ThreePhaseTransformer} & P & P & N & N & P & P & N & N \\
\texttt{ThreePhaseLoad} & P & P & N & N & P & P & N & N \\
\texttt{ThreePhaseGenerator} & P & P & N & N & P & P & N & N \\
\texttt{ThreePhaseExternalGrid} & P & P & N & N & P & P & N & N \\
\bottomrule
\end{longtable}
\normalsize

\subsection{Notable Row-Level Notes}

\begin{itemize}
  \item \texttt{Load} is direct in OPF, but the explicit variable is currently bus-level load shedding in the parity path rather than per-load flexibility.
  \item \texttt{Shunt} is direct in the main PF/OPF path; in DPF it is simplified to a fixed power contribution instead of a true \(V^2\)-dependent admittance.  In RPO, shunts are discrete decision variables (switchable step count).
  \item \texttt{Transformer2W}, \texttt{Transformer3W}, \texttt{Switch}, \texttt{FlexibleLoad}, \texttt{AsymmetricLoad}, and \texttt{Charger} rely on projection and therefore are invisible to DPF and ONR unless the caller has already converted them.
  \item \texttt{RenewableGen} and \texttt{PVSystem} are now parity-IPM decision variables \((P^{\qual{ren}}, Q^{\qual{ren}})\) for curtailable renewables and controllable PV.  In DC OPF they appear as fixed injections.
  \item \texttt{Storage} (AC and DC) and \texttt{DCDCConverter} are now parity-IPM decision variables.  DCDC dispatch includes analytical \(I^2R\) Jacobian and Hessian.
  \item \texttt{FlexibleLoad} is now a parity-IPM decision variable \(P^{\qual{flex}}\),  after projection to canonical form.
  \item DC OPF uses a \(B\)-\(\theta\) network model with linearized or quadratic generator costs and branch flow limits.
  \item RPO uses AC OPF as the inner NLP and treats transformer taps and switchable shunts as discrete variables.
  \item \texttt{VirtualPowerPlant}, \texttt{Microgrid}, and \texttt{MobileStorage} are seen by the main PF/OPF path and by the hybrid overload of DPF, but not by the current AC-only ONR or DC OPF interfaces.
\end{itemize}

\subsection{Interpretation}

The table shows two important facts about the current code base.
\begin{enumerate}
  \item The main PF, native OPF, parity OPF, and SC paths already share a strong canonical-equivalent layer.
  \item The DPF and ONR paths have narrower direct component coverage and should be treated as AC-distribution-specialized modules rather than fully general consumers of every public model struct.
\end{enumerate}
